\documentclass{amsart}
% For dealing with references we use the comment environment
\usepackage{verbatim}
\newenvironment{reference}{\comment}{\endcomment}
%\newenvironment{reference}{}{}
\newenvironment{slogan}{\comment}{\endcomment}
\newenvironment{history}{\comment}{\endcomment}

% For commutative diagrams we use Xy-pic
\usepackage[all]{xy}

% We use 2cell for 2-commutative diagrams.
\xyoption{2cell}
\UseAllTwocells

% We use multicol for the list of chapters between chapters
\usepackage{multicol}

% This is generally recommended for better output
\usepackage{lmodern}
\usepackage[T1]{fontenc}

% For cross-file-references

% Package for hypertext links:
\usepackage{hyperref}

% For any local file, say "hello.tex" you want to link to please

% Theorem environments.
%
\theoremstyle{plain}
\newtheorem{theorem}[subsection]{Theorem}
\newtheorem{proposition}[subsection]{Proposition}
\newtheorem{lemma}[subsection]{Lemma}

\theoremstyle{definition}
\newtheorem{definition}[subsection]{Definition}
\newtheorem{example}[subsection]{Example}
\newtheorem{exercise}[subsection]{Exercise}
\newtheorem{situation}[subsection]{Situation}

\theoremstyle{remark}
\newtheorem{remark}[subsection]{Remark}
\newtheorem{remarks}[subsection]{Remarks}

\numberwithin{equation}{subsection}

% Macros
%
\def\lim{\mathop{\mathrm{lim}}\nolimits}
\def\colim{\mathop{\mathrm{colim}}\nolimits}
\def\Spec{\mathop{\mathrm{Spec}}}
\def\Hom{\mathop{\mathrm{Hom}}\nolimits}
\def\Ext{\mathop{\mathrm{Ext}}\nolimits}
\def\SheafHom{\mathop{\mathcal{H}\!\mathit{om}}\nolimits}
\def\SheafExt{\mathop{\mathcal{E}\!\mathit{xt}}\nolimits}
\def\Sch{\mathit{Sch}}
\def\Mor{\mathop{\mathrm{Mor}}\nolimits}
\def\Ob{\mathop{\mathrm{Ob}}\nolimits}
\def\Sh{\mathop{\mathit{Sh}}\nolimits}
\def\NL{\mathop{N\!L}\nolimits}
\def\CH{\mathop{\mathrm{CH}}\nolimits}
\def\proetale{{pro\text{-}\acute{e}tale}}
\def\etale{{\acute{e}tale}}
\def\QCoh{\mathit{QCoh}}
\def\Ker{\mathop{\mathrm{Ker}}}
\def\Im{\mathop{\mathrm{Im}}}
\def\Coker{\mathop{\mathrm{Coker}}}
\def\Coim{\mathop{\mathrm{Coim}}}

% Boxtimes
%
\DeclareMathSymbol{\boxtimes}{\mathbin}{AMSa}{"02}

%
% Macros for moduli stacks/spaces
%
\def\QCohstack{\mathcal{QC}\!\mathit{oh}}
\def\Cohstack{\mathcal{C}\!\mathit{oh}}
\def\Spacesstack{\mathcal{S}\!\mathit{paces}}
\def\Quotfunctor{\mathrm{Quot}}
\def\Hilbfunctor{\mathrm{Hilb}}
\def\Curvesstack{\mathcal{C}\!\mathit{urves}}
\def\Polarizedstack{\mathcal{P}\!\mathit{olarized}}
\def\Complexesstack{\mathcal{C}\!\mathit{omplexes}}
% \Pic is the operator that assigns to X its picard group, usage \Pic(X)
% \Picardstack_{X/B} denotes the Picard stack of X over B
% \Picardfunctor_{X/B} denotes the Picard functor of X over B
\def\Pic{\mathop{\mathrm{Pic}}\nolimits}
\def\Picardstack{\mathcal{P}\!\mathit{ic}}
\def\Picardfunctor{\mathrm{Pic}}
\def\Deformationcategory{\mathcal{D}\!\mathit{ef}}

\setlength{\emergencystretch}{3em}
\begin{document}
\title[Stacks Project, Tag 0A48]{297 --- Finite etale covers of proper schemes over henselian local rings are determined by the closed fibre}
\maketitle
\noindent Copyright (C) 2005--2025 Johan de Jong. Stacks Project authors. Source: \href{https://stacks.math.columbia.edu/tag/0A48}{Tag 0A48}.

\noindent Editorial difficulty note (not author text). Difficulty H3: One must algebraize the entire compatible infinitesimal tower of a finite etale cover and establish full faithfulness through graph lifting. The proof then uses approximation of solutions over the completion and a directed-limit reduction through henselian Noetherian G-rings to handle an arbitrary henselian local base..

\noindent Editorial provenance note (not author text). History: excerpt from revision \texttt{a0444\allowbreak 6e57e\allowbreak c1fbc\allowbreak 252a8\allowbreak 71afc\allowbreak ec775\allowbreak 2fb28\allowbreak 07b14}, pione.tex, lines 1621--1746. Reference presentation was adapted; original-author context and core support blocks were retained or added. The mathematical wording is unchanged. Licensed under GNU FDL 1.2 or later; no invariant sections or cover texts. See ../licenses/COPYING. Context blocks reproduce author definitions and supporting results; they are not additional counted problems.

\begin{lemma}
\label{lemma-finite-etale-on-proper-over-henselian}
Let $A$ be a henselian local ring. Let $X$ be a proper scheme over $A$
with closed fibre $X_0$. Then the functor
$$
\textit{F\'Et}_X \to \textit{F\'Et}_{X_0},\quad
U \longmapsto U_0 = U \times_X X_0
$$
is an equivalence of categories.
\end{lemma}

\begin{proof}
The proof given here is an example of applying algebraization and
approximation. We proceed in a number of stages.

\medskip\noindent
Essential surjectivity when $A$ is a complete local Noetherian ring.
Let $X_n = X \times_{\Spec(A)} \Spec(A/\mathfrak m^{n + 1})$.
By \'Etale Morphisms, Theorem \href{https://stacks.math.columbia.edu/tag/039R}{theorem remarkable equivalence (Tag 039R)}
the inclusions
$$
X_0 \to X_1 \to X_2 \to \ldots
$$
induce equivalence of categories between the category
of schemes \'etale over $X_0$ and the category of schemes
\'etale over $X_n$.
Moreover, if $U_n \to X_n$ corresponds to a finite \'etale
morphism $U_0 \to X_0$, then $U_n \to X_n$ is finite too, for example
by More on Morphisms, Lemma
\href{https://stacks.math.columbia.edu/tag/09ZW}{lemma thicken property morphisms cartesian (Tag 09ZW)}.
In this case the morphism $U_0 \to \Spec(A/\mathfrak m)$
is proper as $X_0$ is proper over $A/\mathfrak m$. Thus we may apply
Grothendieck's algebraization theorem
(in the form of
Cohomology of Schemes, Lemma
\href{https://stacks.math.columbia.edu/tag/09ZT}{lemma algebraize formal scheme finite over proper (Tag 09ZT)})
to see that there is a finite morphism $U \to X$ whose restriction
to $X_0$ recovers $U_0$. By More on Morphisms, Lemma
\href{https://stacks.math.columbia.edu/tag/0A43}{lemma check smoothness on infinitesimal nbhds (Tag 0A43)}
we see that $U \to X$ is \'etale at every point of $U_0$.
However, since every point of $U$ specializes to a point of $U_0$
(as $U$ is proper over $A$), we conclude that $U \to X$ is \'etale.
In this way we conclude the functor is essentially surjective.

\medskip\noindent
Fully faithfulness when $A$ is a complete local Noetherian ring.
Let $U \to X$ and $V \to X$ be finite \'etale morphisms and
let $\varphi_0 : U_0 \to V_0$ be a morphism over $X_0$. Look at
the morphism
$$
\Gamma_{\varphi_0} : U_0 \longrightarrow U_0 \times_{X_0} V_0
$$
This morphism is both finite \'etale and a closed immersion.
By essential surjectivity applied to $X = U \times_X V$ we find
a finite \'etale morphism $W \to U \times_X V$ whose special
fibre is isomorphic to $\Gamma_{\varphi_0}$. Consider the projection
$W \to U$. It is finite \'etale and an isomorphism over $U_0$ by
construction. By \'Etale Morphisms, Lemma
\href{https://stacks.math.columbia.edu/tag/04DH}{lemma finite etale one point (Tag 04DH)}
$W \to U$ is an isomorphism in an open neighbourhood of $U_0$.
Thus it is an isomorphism and the composition $\varphi : U \cong W \to V$
is the desired lift of $\varphi_0$.

\medskip\noindent
Essential surjectivity when $A$ is a henselian local Noetherian G-ring.
Let $U_0 \to X_0$ be a finite \'etale morphism.
Let $A^\wedge$ be the completion of $A$ with respect to the maximal ideal.
Let $X^\wedge$ be the base change of $X$ to $A^\wedge$.
By the result above there exists a finite \'etale morphism
$V \to X^\wedge$ whose special fibre is $U_0$.
Write $A^\wedge = \colim A_i$ with $A \to A_i$ of finite type.
By Limits, Lemma \href{https://stacks.math.columbia.edu/tag/01ZM}{lemma descend finite presentation (Tag 01ZM)}
there exists an $i$ and a finitely presented morphism $U_i \to X_{A_i}$
whose base change to $X^\wedge$ is $V$. After increasing $i$
we may assume that $U_i \to X_{A_i}$ is finite and \'etale
(Limits, Lemmas \href{https://stacks.math.columbia.edu/tag/01ZO}{lemma descend finite finite presentation (Tag 01ZO)} and
\href{https://stacks.math.columbia.edu/tag/07RP}{lemma descend etale (Tag 07RP)}). Writing
$$
A_i = A[x_1, \ldots, x_n]/(f_1, \ldots, f_m)
$$
the ring map $A_i \to A^\wedge$ can be reinterpreted as a solution
$(a_1, \ldots, a_n)$ in $A^\wedge$ for the system of equations $f_j = 0$.
By Smoothing Ring Maps, Theorem \href{https://stacks.math.columbia.edu/tag/07QY}{theorem approximation property (Tag 07QY)}
we can approximate this solution (to order $11$ for example) by a solution
$(b_1, \ldots, b_n)$ in $A$. Translating back we find an $A$-algebra map
$A_i \to A$ which gives the same closed point as the original map
$A_i \to A^\wedge$ (as $11 > 1$). The base change $U \to X$ of $V \to X_{A_i}$
by this ring map will therefore be a finite \'etale morphism whose
special fibre is isomorphic to $U_0$.

\medskip\noindent
Fully faithfulness when $A$ is a henselian local Noetherian G-ring.
This can be deduced from essential surjectivity in exactly the same
manner as was done in the case that $A$ is complete Noetherian.

\medskip\noindent
General case. Let $(A, \mathfrak m)$ be a henselian local ring.
Set $S = \Spec(A)$ and denote $s \in S$ the closed point. By Limits, Lemma
\href{https://stacks.math.columbia.edu/tag/0A0P}{lemma proper limit of proper finite presentation noetherian (Tag 0A0P)}
we can write $X \to \Spec(A)$ as a cofiltered limit of
proper morphisms $X_i \to S_i$ with $S_i$ of finite type over $\mathbf{Z}$.
For each $i$ let $s_i \in S_i$ be the image of $s$.
Since $S = \lim S_i$ and $A = \mathcal{O}_{S, s}$ we have
$A = \colim \mathcal{O}_{S_i, s_i}$. The ring $A_i = \mathcal{O}_{S_i, s_i}$
is a Noetherian local G-ring (More on Algebra, Proposition
\href{https://stacks.math.columbia.edu/tag/07PX}{proposition ubiquity G ring (Tag 07PX)}).
By More on Algebra, Lemma \href{https://stacks.math.columbia.edu/tag/0A04}{lemma henselization colimit (Tag 0A04)}
we see that $A = \colim A_i^h$. By
More on Algebra, Lemma \href{https://stacks.math.columbia.edu/tag/07QR}{lemma henselization G ring (Tag 07QR)}
the rings $A_i^h$ are G-rings. Thus we see that $A = \colim A_i^h$ and
$$
X = \lim (X_i \times_{S_i} \Spec(A_i^h))
$$
as schemes. The category of schemes finite \'etale over $X$ is the limit
of the category of schemes finite \'etale over
$X_i \times_{S_i} \Spec(A_i^h)$ (by
Limits, Lemmas
\href{https://stacks.math.columbia.edu/tag/01ZM}{lemma descend finite presentation (Tag 01ZM)},
\href{https://stacks.math.columbia.edu/tag/01ZO}{lemma descend finite finite presentation (Tag 01ZO)}, and
\href{https://stacks.math.columbia.edu/tag/07RP}{lemma descend etale (Tag 07RP)})
The same thing is true for schemes finite \'etale over
$X_0 = \lim (X_i \times_{S_i} s_i)$.
Thus we formally deduce the result for $X / \Spec(A)$
from the result for the $(X_i \times_{S_i} \Spec(A_i^h)) / \Spec(A_i^h)$
which we dealt with above.
\end{proof}
\end{document}
